% الجزء: sets-functions-relations
% الفصل: sets
% القسم: russells-paradox

\documentclass[../../../include/open-logic-section]{subfiles}


\begin{document}

\olfileid[ar]{sfr}{set}{rus}
\olsection{مفارقة راسل}

سوّغت الامتدادية أن نرمز بـ$\Setabs{\OLId{latin-ordinary}{x}}{\phi(\OLId{latin-ordinary}{x})}$ إلى \emph{مجموعة}
قيم $\OLId{latin-ordinary}{x}$ التي تحقق~$\phi(\OLId{latin-ordinary}{x})$. ولكن معناها \emph{على التحقيق} شرط:
\emph{إذا} وجدت مجموعة تضم كل ما يتصف بـ~$\phi$ ولا تضم غيره،
\emph{فإنها} واحدة. فمتى ثبتت الخاصية~$\phi$ كانت
$\Setabs{\OLId{latin-ordinary}{x}}{\phi(\OLId{latin-ordinary}{x})}$ فريدة \emph{إن وجدت}.

ولا غنى عن هذا الشرط!{} فليس كل وصف صالحًا لـ\emph{الاستيعاب}؛
بل من الخواص ما \emph{لا} يعيّن مجموعة. ولو جاز ذلك لكل خاصية لزم
التناقض الصريح، وأشهر شاهد عليه مفارقة راسل.

قد تكون المجموعات !!{element}s لمجموعات أخرى؛ فمجموعة قوى~$\OLId{latin-looped}{A}$ مؤلفة
من مجموعات. فلا حرج إذن في السؤال: هل هذه المجموعة !!a{element}
من تلك؟ ثم نسأل: أيجوز أن تكون المجموعة عضوًا في نفسها؟
ليس في مجرد تصور المجموعة ما يبدو مانعًا منه. فقد يقال إن
\emph{جميع} المجموعات أشياء مجتمعة، فيجوز أن تضمها مجموعة واحدة
تسمى مجموعة جميع المجموعات. وهذه، لو وجدت، لكانت مجموعة، ومن ثم
كانت !!a{element} من مجموعة جميع المجموعات.

وأما مفارقة راسل فتنشأ من وصف المجموعة بأنها ليست !!a{element} من
نفسها، أي بأنها \emph{غير منتمية إلى نفسها}. فلنفرض أن لكل
المجموعات التي ليست !!a{element} من أنفسها مجموعة جامعة:
\[
\OLId{latin-looped}{R} = \Setabs{\OLId{latin-ordinary}{x}}{\OLId{latin-ordinary}{x} \notin \OLId{latin-ordinary}{x}}
\]
أيمكن وجودها؟ سنبرهن على امتناعه.

\begin{thm}[مفارقة راسل]\ollabel{thm:russells-paradox}
يمتنع وجود مجموعة $\OLId{latin-looped}{R} = \Setabs{\OLId{latin-ordinary}{x}}{\OLId{latin-ordinary}{x} \notin \OLId{latin-ordinary}{x}}$.
\end{thm}

\begin{proof}
لنفرض وجود $\OLId{latin-looped}{R} = \Setabs{\OLId{latin-ordinary}{x}}{\OLId{latin-ordinary}{x} \notin \OLId{latin-ordinary}{x}}$. فيلزم من تعريفها
$\OLId{latin-looped}{R} \in \OLId{latin-looped}{R}$ إذا وفقط إذا كان $\OLId{latin-looped}{R} \notin \OLId{latin-looped}{R}$؛ وهذا تناقض.
\end{proof}

\begin{tagblock}{novice}
\begin{explain}
وتفصيل البرهان أن~$\OLId{latin-looped}{R}$، إن وجدت، صح السؤال عن $\OLId{latin-looped}{R} \in \OLId{latin-looped}{R}$.
فإن فرضنا $\OLId{latin-looped}{R} \in \OLId{latin-looped}{R}$، تذكرنا أن~$\OLId{latin-looped}{R}$ إنما عرّفت بجمع المجموعات التي
ليست !!{element}s في أنفسها. فمع $\OLId{latin-looped}{R} \in \OLId{latin-looped}{R}$ لا تتصف~$\OLId{latin-looped}{R}$ بالخاصية
التي تعرّف~$\OLId{latin-looped}{R}$. وليس في~$\OLId{latin-looped}{R}$ إلا ما اتصف بها؛ فلا تكون~$\OLId{latin-looped}{R}$
!!a{element} من~$\OLId{latin-looped}{R}$، أي $\OLId{latin-looped}{R} \notin \OLId{latin-looped}{R}$. وهذا تناقض، إذ يستحيل
أن تكون~$\OLId{latin-looped}{R}$ !!a{element} من~$\OLId{latin-looped}{R}$ وألا تكون كذلك في آن واحد.

فقد بطل فرض $\OLId{latin-looped}{R} \in \OLId{latin-looped}{R}$، وثبت $\OLId{latin-looped}{R} \notin \OLId{latin-looped}{R}$. غير أن هذا أيضًا
يفضي إلى التناقض!{} إذ متى كان $\OLId{latin-looped}{R} \notin \OLId{latin-looped}{R}$ اتصفت~$\OLId{latin-looped}{R}$ نفسها
بالخاصية المعرّفة لـ~$\OLId{latin-looped}{R}$، فصارت~$\OLId{latin-looped}{R}$ !!a{element} من~$\OLId{latin-looped}{R}$، كسائر
المجموعات التي لا تنتمي إلى أنفسها. فعاد المحال: تكون !!a{element}
من~$\OLId{latin-looped}{R}$ ولا تنتمي إليها معًا.
\end{explain}
\end{tagblock}

\begin{digress}
فكيف توضع نظرية مجموعات تسلم من مفارقة راسل، فلا تقضي بالقول
\emph{غير المتسق} إن $\OLId{latin-looped}{R} = \Setabs{\OLId{latin-ordinary}{x}}{\OLId{latin-ordinary}{x} \notin \OLId{latin-ordinary}{x}}$ موجودة؟
سبيل ذلك بديهيات تضبط بدقة شروط وجود المجموعات وعدم وجودها.

وليست النظرية الموجزة في هذا الفصل من ذلك القبيل؛ فهي
\emph{ساذجة حقًا}. % حفظ أمر تشكيل غير منفذ من المصدر: \"
وإنما تقضي بالامتدادية، وتجيز، إذا أعطيت مجموعات، تكوين اتحادها
وتقاطعها ونحوهما. وللنظرية بناء أشد إحكامًا يمكن سلوكه.
\oliflabeldef{cumul:::part}{وسيكون محل هذا الإحكام الجزء
\olref[cumul][][]{part}. أما هنا فنمضي على الطريقة الساذجة مع التحفظ
في اجتناب التناقض. % حفظ أمر تشكيل غير منفذ من المصدر: \"
}{}
\end{digress}


\end{document}
